\documentclass[10pt,a4paper]{article}

\usepackage{fontspec}
\usepackage[a4paper,top=1.1cm,bottom=1.1cm,left=1.5cm,right=1.5cm]{geometry}
\usepackage{xcolor}
\usepackage{titlesec}
\usepackage{enumitem}
\usepackage{hyperref}
\usepackage{microtype}
\usepackage{parskip}
\usepackage{ragged2e}
\usepackage{needspace}

% Fonts — sans-serif throughout. Loaded by filename so kpsewhich resolves
% them from the TeX tree on any platform (avoids relying on OS font registry).
\setmainfont{texgyreheros-regular.otf}[
  ItalicFont    = texgyreheros-italic.otf,
  BoldFont      = texgyreheros-bold.otf,
  BoldItalicFont= texgyreheros-bolditalic.otf,
  Ligatures=TeX,
  Scale=0.92
]
\newfontfamily\headerfont{texgyreheros-regular.otf}[
  ItalicFont    = texgyreheros-italic.otf,
  BoldFont      = texgyreheros-bold.otf,
  BoldItalicFont= texgyreheros-bolditalic.otf,
  Ligatures=TeX
]

% Colors
\definecolor{accent}{HTML}{1F3A5F}
\definecolor{subtle}{HTML}{555555}
\definecolor{rule}{HTML}{B0B0B0}

% Hyperlinks
\hypersetup{
  colorlinks=true,
  urlcolor=accent,
  linkcolor=accent
}

% Section formatting
\titleformat{\section}
  {\headerfont\large\bfseries\color{accent}}
  {}{0em}
  {\MakeUppercase}
  [\vspace{-0.4em}{\color{rule}\rule{\linewidth}{0.5pt}}]
\titlespacing*{\section}{0pt}{0.6em}{0.3em}

% Subsection (job entries)
\newcommand{\role}[4]{%
  \needspace{4\baselineskip}%
  \noindent\textbf{#1} \textemdash{} \textit{#2} \hfill {\small\color{subtle}#3}\\[-0.25em]
  {\small\color{subtle}#4}\par\vspace{0.1em}
}

% Compact list
\setlist[itemize]{
  leftmargin=1.1em,
  topsep=0.1em,
  itemsep=0.1em,
  parsep=0pt,
  partopsep=0pt
}

\pagenumbering{gobble}
\emergencystretch=2em

\begin{document}

% =============== HEADER ===============
\noindent
{\headerfont\Huge\bfseries\color{accent} Iaroslav Postovalov}\\[0.4em]
{\footnotesize
  Hannover, Germany \;\textbar\; \href{mailto:postovalovya@gmail.com}{postovalovya@gmail.com} \;\textbar\; \href{https://linkedin.com/in/iaroslav-postovalov}{linkedin.com/in/iaroslav-postovalov} \;\textbar\; \href{https://github.com/commandertvis}{github.com/commandertvis} \;\textbar\; \href{https://commandertvis.github.io}{commandertvis.github.io}
}

\vspace{0.3em}

% =============== SUMMARY ===============
\noindent
Compiler and JVM developer, 21 years old. Sole maintainer of a production blockchain language (\textasciitilde197K LOC Kotlin); shipped a new immutable serializable IR and a GraalVM Truffle peer backend alongside the existing tree-walking interpreter. On the Kotlin/JVM compiler team at JetBrains from age 17, part-time through university: IR backend, bytecode optimizations, inline class interop. Published research on JIT compilation targeting JVM, JavaScript, and WebAssembly.

\vspace{0.3em}
\noindent
\textbf{Languages:} native Russian, fluent English \quad\textbar\quad \textbf{Legal status:} Russian citizenship, EU Blue Card (Germany)

% =============== EXPERIENCE ===============
\section*{Experience}

\role{Software Engineer}{ChromaWay}{Apr 2025 -- Present}{Remote \textbullet{} Full-time (concurrent with online education)}
\begin{itemize}
  \item Sole maintainer (since Feb 2026) of Rell (\href{https://github.com/ChromiaProject/rell}{github.com/ChromiaProject/rell}), the smart-contract language of the Chromia relational blockchain. \textasciitilde197K LOC Kotlin; tree-walking AST interpreter, operators lowered to SQL against PostgreSQL.
  \item New compiler IR \textemdash{} immutable, fully-resolved tree replacing a runtime-coupled AST. FlatBuffers wire format with bit-identical round-trip gate per CI build; drove the module split into separate frontend, IR, and runtime.
  \item Second backend on the IR: a self-specializing AST interpreter on GraalVM Truffle (partial-evaluation JIT framework). CI gates for identical results across tree-walker, IR round-trip, and Truffle.
  \item Replaced a hand-written combinator parser with an ANTLR4 grammar. Rebuilt PostgreSQL emission on jOOQ (typed SQL DSL), eliminating pre-existing bugs.
  \item Co-authored runtime state snapshots for chains. PostgreSQL: query-cancellation propagation, per-process schema isolation for parallel tests, removal of a catalog-lock deadlock under concurrent chain boot.
  \item New language features: \texttt{struct.copy()}, \texttt{try\_call\_catch} (structured error handling with automatic database rollback).
  \item Maintainer of the LSP server, client stub generator, and stdlib doc generator; IntelliJ plugin with Chromia CLI integration. Maven$\rightarrow$Gradle and JUnit 4$\rightarrow$5 migrations.
  \item Performance harness: JMH (Java Microbenchmark Harness) suites, Async Profiler automation. Per-commit performance reports on GitHub Pages. External-project regression toolkit in CI.
  \item Codebase health: dependency cuts to stdlib (Guava, Apache Commons, HttpClient; Picocli $\rightarrow$ Clikt); locale-determinism fixes; Python integration tests rewritten in Kotlin.
\end{itemize}

\role{Software Developer}{JetBrains (Kotlin Team)}{Sep 2022 -- Dec 2024}{Yerevan, then Bremen (Germany) \textbullet{} Part-time (concurrent with full-time education)}
\begin{itemize}
  \item Kotlin/JVM compiler team during the K2 frontend transition. IR backend, performance work, language features.
  \item Memory allocation reductions in compiler frontend and IR lowering passes. Bytecode-level optimizations targeting boxing conversions in generated code.
  \item \texttt{@JvmExposeBoxed} annotation: seamless integration of inline/value classes with Java callers.
  \item Android testing infrastructure migration to the official Android Gradle Plugin framework. Compiler test renovation for the new K2 frontend, including test runs on Android.
  \item Bug fixes in IR serialization and lazy declaration deserialization causing production compiler crashes.
\end{itemize}

\role{Software Developer (Intern)}{MiLaboratories}{Jul -- Aug 2022}{Remote}
\begin{itemize}
  \item Kotlin DSL with statistical and serialization features for an internal plotting library based on \texttt{lets-plot-kotlin}.
\end{itemize}

\role{Researcher}{JetBrains Research, Nuclear Physics Methods Lab}{Sep 2020 -- Apr 2022}{Dolgoprudny (Russia), then Remote}
\begin{itemize}
  \item Group's research scope: non-accelerator particle physics, numerical simulations, and software development for experimental physics.
  \item Core developer of KMath (\href{https://github.com/SciProgCentre/kmath}{github.com/SciProgCentre/kmath}, 700+ stars), a cross-platform mathematical library for Kotlin. Core design contributions to the abstract algebra API, separating operations from data structures. N-dimensional array support, automatic differentiation, integrations with math libraries (EJML, ND4J, Commons Math).
  \item Cross-platform expression compiler targeting JVM bytecode (ObjectWeb ASM), JavaScript, and WebAssembly from a single AST. Boxing/unboxing optimization pass with bytecode type state tracking. Paper: ``Compilation of mathematical expressions in Kotlin'' (arXiv:2102.07924).
  \item \texttt{kmath-gsl}: Kotlin/Native bindings to the GNU Scientific Library via C interop, implementing matrices and vectors. Native build/toolchain setup, CI, documentation, performance benchmarking.
  \item KMath-to-LaTeX renderer for computable expressions in Jupyter notebooks.
\end{itemize}

\role{Research Intern}{Lab. Math. Modeling, Lavrentiev Lyceum No. 130}{Sep 2019 -- Apr 2021}{Novosibirsk (Russia)}
\begin{itemize}
  \item Interactive system for studying probability distributions used in Monte Carlo methods (\href{https://gitlab.com/CMDR_Tvis/monte-carlo-script}{gitlab.com/CMDR\_Tvis/monte-carlo-script}). Custom DSL for defining formulas, with C++ tree interpreter and bytecode interpreter for performance research.
  \item Kotlin/Ktor REST API for plot generation and function evaluation. Cross-platform bindings: Python via \texttt{ctypes}, Java/Kotlin via JNI. Diploma at the 58th International Scientific Student Conference (MNSK-2020).
\end{itemize}

% =============== TECHNICAL SKILLS ===============
\needspace{6\baselineskip}
\section*{Technical Skills}

\noindent
\small
\begin{tabular}{@{}p{3.4cm}@{\hspace{0.5em}}p{14cm}@{}}
\textbf{Languages} & Kotlin (expert), Java, C, Python, C\#, SQL \\[0.2em]
\textbf{Compiler engineering} & Language design; lexer/parser implementation (ANTLR4, SLL/LL with error recovery); AST and IR design; immutable serializable IR (FlatBuffers); lowering passes; type inference and checking; semantic analysis; code generation; optimization passes (constant folding, inlining, boxing/unboxing elimination); tree-walk and AST-rewriting interpreters; multi-target backends (JVM bytecode, JavaScript, WebAssembly, SQL); differential and end-to-end compiler testing \\[0.2em]
\textbf{JVM internals} & Bytecode generation (ObjectWeb ASM); invokedynamic; inline/value classes; allocation profiling; class loading; JNI/FFI; GraalVM; Truffle language framework; partial evaluation; polyglot embedding \\[0.2em]
\textbf{Performance} & Microbenchmarking (JMH); profiling (JFR, async-profiler); performance regression analysis; test automation \\[0.2em]
\textbf{Tooling \& infra} & Git, Gradle, GitLab CI, PostgreSQL, jOOQ, Docker, JUnit 5, Linux, Kotlin Multiplatform, open source maintenance \\[0.2em]
\textbf{Domain} & Blockchain smart contract languages, SQL codegen, scientific computing, mathematical library design, automatic differentiation, AI-assisted development \\
\end{tabular}
\normalsize

% =============== EDUCATION ===============
\section*{Education}

\role{BSc Applied Computer Science}{Constructor University Bremen}{Sep 2023 -- Aug 2027}{}

\role{Applied Mathematics \& CS}{Lomonosov Moscow State University, Yerevan Branch}{Sep 2022 -- Jun 2023}{}

% =============== PUBLICATIONS & HONORS ===============
\section*{Publications \& Honors}

\begin{itemize}
  \item \textbf{Preprint:} \textit{Compilation of mathematical expressions in Kotlin} --- arXiv, Feb 2021 (\href{https://arxiv.org/abs/2102.07924}{arXiv:2102.07924}).
  \item \textbf{Paper:} \textit{NMPUD: a computer system for sampling and examining probability distributions} --- ITMM-2020 proceedings, Tomsk, 2021. Co-authored with Prof. A. Voytishek.
  \item \textbf{Talk:} \textit{Dynamic compilation of mathematical expressions with Kotlin} --- SnowOne 2022 (JUGNsk), Novosibirsk.
  \item \textbf{Winner, National Technology Olympiad} --- ``Automation of Business Processes'' track (1C partner case). HSE University, Moscow \textbullet{} Mar 2022.
  \item \textbf{Volunteering:} Translator at Translators without Borders (Apr 2022 -- Feb 2023, 11 months).
\end{itemize}

\end{document}
